\documentclass[10pt]{article}
\usepackage[margin=0.75in]{geometry}
\usepackage{booktabs}
\usepackage{graphicx}
\usepackage{hyperref}

\title{Communication-Aware CUDA and MPI Training for a Mini-Transformer}
\author{}
\date{}

\begin{document}
\maketitle

\section{Introduction and Motivation}
Training modern neural networks is usually presented through high-level
frameworks, but the performance-critical work underneath is a tightly coupled
systems problem.  Each training step mixes dense matrix multiplication,
memory-bound elementwise kernels, reductions, optimizer updates, and, in the
distributed case, gradient communication.  The goal of this project is to make
those interactions visible by building and optimizing a mini-Transformer
training system directly in CUDA and MPI.  The model is intentionally small
enough to understand end to end, but it contains the same bottleneck classes as
larger production systems: attention, feed-forward GEMMs, LayerNorm, softmax,
embedding lookups, Adam, and data-parallel allreduce.

Our starting point was a fully custom CUDA implementation with MPI data
parallelism.  This baseline was useful because every operation had a clear
owner: if a kernel was slow, we could inspect its memory traffic, launch
geometry, and synchronization behavior directly.  The first major optimization
was to replace the largest GEMMs with cuBLAS and then with Tensor Core GEMMs
using FP32 inputs and FP32 accumulation with fast FP16 internal math.  This
changed the performance profile substantially.  At larger hidden sizes the
GEMMs became much faster, while the remaining LayerNorm, activation, softmax,
and Adam kernels became a larger fraction of step time.  This is the central
theme of the project: each successful optimization changes the bottleneck, so
the next decision must be made from fresh measurements rather than from a fixed
intuition about what should be slow.

The distributed results add another layer.  On four GPUs, gradient
synchronization can dominate the small models even when the single-GPU kernels
are efficient.  We therefore evaluate blocking allreduce, bucketed overlap,
backend selection, and communication noise carefully.  One surprising result is
that the best H128 backend differs between single-GPU and four-GPU execution:
Tensor Core cuBLAS wins on one GPU, while the custom CUDA kernels win in the
four-GPU MPI setting.  Profiling shows that the Tensor Core path reduces GPU
kernel time, but for the small per-rank H128 distributed workload its fixed
runtime and synchronization overheads outweigh the compute advantage.  This
kind of crossover is exactly why the project emphasizes measurement rather
than treating any library or fusion strategy as automatically better.

The final objective is not only to make the code faster, but to explain why
each optimization worked, failed, or remained inconclusive.  Negative results
are part of that story: cuBLASLt epilogue fusion was slower for these shapes,
CUDA Graphs provided limited benefit while the GPU was already busy, and some
MPI overlap strategies were limited by host staging and variance.  By combining
correctness tests, repeated benchmarks, Nsight Systems and Nsight Compute
profiles, and a roofline analysis, this report presents a performance
engineering narrative for CUDA and MPI training that is small enough to audit
but rich enough to expose real architectural tradeoffs.

\section{System Architecture}
% May 15 target: 500 words.

\section{Correctness Methodology}
% May 16 target: 500 words.

\section{Benchmarking Methodology}
% May 17 target: 500 words.

\section{Single-GPU Optimization Results}

\section{Distributed Optimization Results}

\section{Roofline and Bottleneck Analysis}

\section{Discussion and Limitations}

\section{Conclusion}

\end{document}
